\chapter{Conclusão}
\label{cap:conclusao}

\section{Síntese}

Este trabalho demonstrou que é possível eliminar inteiramente a etapa de
reamostragem em grade latitude-longitude regular do processo de geração
de condição inicial e de fronteira lateral do \mpasa{}, interpolando
diretamente entre a malha de Voronoi nativa de origem (uma rodada global
já concluída) e a malha de Voronoi nativa de destino (o domínio
regional), e escrevendo os arquivos \texttt{init.nc}/\texttt{lbc.*.nc}
consumidos pelo núcleo dinâmico \texttt{mpas\_atmosphere} sem depender do
\texttt{init\_atmosphere\_model} original. A viabilidade dessa rota não
era óbvia \emph{a priori}: a investigação do código-fonte de referência
mostrou que o formato binário intermediário do WPS é, por definição, uma
grade regular, de modo que nenhum ``escopo mínimo'' preservaria o ganho
de precisão da interpolação nativa sem substituir inteiramente o
caminho -- horizontal, vertical, hidrostático e de superfície -- por
implementação própria.

O método central -- interpolação baricêntrica na malha dual de Delaunay
-- não é uma proposta original deste trabalho, mas a aplicação de uma
técnica já validada pela comunidade (via o sistema de assimilação
MPAS-DART) a um novo contexto de uso. A contribuição prática está na
extração cuidadosa, tradução literal e validação exaustiva de cada
fórmula fisicamente necessária do código-fonte de referência do
MPAS-Model -- um processo de engenharia reversa disciplinado que revelou,
ao longo do caminho, tanto suposições incorretas do planejamento inicial
quanto bugs reais de implementação, todos documentados com sua
investigação e correção nos capítulos técnicos deste relatório.

Um padrão metodológico se repete em praticamente todos os problemas mais
difíceis encontrados: a causa raiz raramente estava onde o sintoma
aparecia. A divergência de $\sim\SI{2}{\kilo\metre}$ na grade vertical
não estava no algoritmo de suavização (onde o sintoma se manifestava),
mas em uma suposição incorreta sobre um valor padrão de configuração
calculado dezenas de linhas antes. Os dois bugs que impediam a execução
real do modelo não eram de física nem de numérica -- eram de formato de
uma única variável de metadado (\texttt{xtime}) inicialmente considerada
secundária. Em ambos os casos, uma validação aparentemente completa
(revisão de código linha a linha; comparação numérica campo a campo)
não foi suficiente até que a validação fosse estendida às dependências
reais do problema -- o código \emph{upstream} do trecho suspeito, no
primeiro caso; o consumidor final de verdade, no segundo.

\section{Estado atual e alcance da validação}

Ao final deste trabalho, a rota nativa \emph{voronoi-to-voronoi} está
validada em duas camadas independentes: numericamente, campo a campo,
contra arquivos \texttt{init.nc}/\texttt{lbc.*.nc} reais de produção; e
funcionalmente, executando o \texttt{mpas\_atmosphere} real a partir
desses arquivos e obtendo uma previsão de 24 horas fisicamente sã, cuja
divergência frente à rota de referência cresce de forma suave e
limitada, consistente com o comportamento esperado de duas trajetórias
vizinhas de um sistema dinâmico caótico -- não com um artefato de
implementação. Essa validação, no entanto, foi conduzida em um único
par malha/caso de estudo (\cref{cap:discussao}); a rota foi generalizada
estruturalmente (leitura de \textit{namelist} real em tempo de execução)
para não depender de recompilação em outros cenários, mas ainda não foi
exercitada neles.

\section{Trabalhos futuros}

Três direções concretas de continuidade emergem diretamente das
limitações documentadas no \cref{cap:discussao}:

\begin{enumerate}
  \item \textbf{Implementar a mistura de terreno de fronteira}
    (\texttt{config\_blend\_bdy\_terrain}), fechando o resíduo residual
    de $\sim\SI{171}{\metre}$ concentrado no anel de células de
    fronteira -- a lacuna de maior impacto numérico mensurável
    remanescente, e a mais diretamente relevante caso a fronteira
    lateral se torne um ponto crítico de precisão em usos futuros.
  \item \textbf{Ampliar a validação para outras malhas, resoluções e
    estações do ano}, exercitando a generalização estrutural já
    implementada (\cref{sec:generalizacao}) e verificando se os 61
    níveis de pressão fixa intermediários -- derivados estatisticamente
    de uma única malha/caso -- permanecem adequados em contextos com
    perfil vertical de pressão substancialmente distinto.
  \item \textbf{Avaliar interpolantes de maior continuidade} sobre a
    malha dual de Delaunay (por exemplo, a família de interpolantes
    generalizados de Sibson/Laplace descrita por Hiyoshi e Sugihara
    \cite{hiyoshi2000continuity}) como substituto opcional da
    interpolação baricêntrica linear, caso alguma aplicação futura do
    mesmo motor de remapeamento -- por exemplo, fora do contexto de
    condição inicial/fronteira, em algum uso que dependa diretamente de
    gradiente horizontal suave -- se mostre sensível à descontinuidade
    de derivada nas arestas dos triângulos de Delaunay.
\end{enumerate}

Adicionalmente, uma investigação independente e de menor prioridade
identificada ao longo deste trabalho, mas fora de seu escopo direto: o
campo \texttt{lbc\_qv} dos arquivos de referência reais de produção
mostrou-se constante em todos os níveis verticais por célula
(\cref{cap:resultados}) -- um comportamento fisicamente implausível que
sugere um problema já existente na cadeia de geração da rota original
(via WPS), não introduzido por este trabalho, e que mereceria
investigação própria caso a fronteira lateral de referência continue
sendo usada como padrão de comparação em trabalhos futuros.
